% method.tex
%
% Method section for OffTrackGS. Intended to be \input{} (or \subfile{}) from
% the paper's main.tex. Assumes the main document preamble already loads
% amsmath, amssymb (and, if figures are added, graphicx).
%
%   % method.tex
%
% Method section for OffTrackGS. Intended to be \input{} (or \subfile{}) from
% the paper's main.tex. Assumes the main document preamble already loads
% amsmath, amssymb (and, if figures are added, graphicx).
%
%   % method.tex
%
% Method section for OffTrackGS. Intended to be \input{} (or \subfile{}) from
% the paper's main.tex. Assumes the main document preamble already loads
% amsmath, amssymb (and, if figures are added, graphicx).
%
%   % method.tex
%
% Method section for OffTrackGS. Intended to be \input{} (or \subfile{}) from
% the paper's main.tex. Assumes the main document preamble already loads
% amsmath, amssymb (and, if figures are added, graphicx).
%
%   \input{tex/method}
%
% Notation used throughout this section:
%   scalars      x
%   vectors      \mathbf{x}
%   matrices     \mathbf{X}
%   sets/scenes  \mathcal{X}

\section{Method}
\label{sec:method}

This section describes the proposed method in three parts. Section
\ref{sec:method:3dgs} presents the deformable 3D Gaussian Splatting
representation on which our approach is built. Section
\ref{sec:method:gcsr} introduces Geometry-Constrained Sideview
Regularization (GCSR), a training-time regularizer that constrains scene
geometry at viewpoints not present in the training trajectory. Section
\ref{sec:method:rtvc} introduces Round-Trip View Consistency (RTVC), an
evaluation protocol that quantifies reconstruction quality at such
viewpoints without requiring ground-truth imagery to be captured there.

\subsection{Deformable 3D Gaussian Splatting}
\label{sec:method:3dgs}

\paragraph{Canonical representation.}
The reconstructed tissue is represented as a set of $N$ canonical 3D
Gaussians $\mathcal{G} = \{(\boldsymbol{\mu}_i, \mathbf{s}_i, \mathbf{q}_i,
\alpha_i, \mathbf{c}_i)\}_{i=1}^{N}$, where $\boldsymbol{\mu}_i \in
\mathbb{R}^3$ denotes the mean, $\mathbf{s}_i \in \mathbb{R}^3_{>0}$ the
per-axis scale, $\mathbf{q}_i$ a unit quaternion specifying orientation,
$\alpha_i \in [0,1]$ an opacity, and $\mathbf{c}_i$ a set of
spherical-harmonics coefficients encoding view-dependent color. The
covariance of Gaussian $i$ is given by
\begin{equation}
\boldsymbol{\Sigma}_i = \mathbf{R}(\mathbf{q}_i)\,\mathrm{diag}(\mathbf{s}_i)^2\,\mathbf{R}(\mathbf{q}_i)^\top ,
\end{equation}
where $\mathbf{R}(\mathbf{q}_i)$ is the rotation matrix corresponding to
$\mathbf{q}_i$.

\paragraph{Deformation network.}
Because surgical tissue deforms over the course of a video, $\mathcal{G}$ is
not treated as static. A deformation network $F_\theta$ maps each canonical
Gaussian and a query time $t$ to a set of residuals,
\begin{equation}
\big(\Delta\boldsymbol{\mu}_i(t),\, \Delta\mathbf{s}_i(t),\, \Delta\mathbf{q}_i(t),\, \Delta\alpha_i(t)\big) = F_\theta(\boldsymbol{\mu}_i, t),
\end{equation}
which define the time-$t$ Gaussian $\boldsymbol{\mu}_i(t) =
\boldsymbol{\mu}_i + \Delta\boldsymbol{\mu}_i(t)$, and analogously for
$\mathbf{s}_i(t)$, $\mathbf{q}_i(t)$, and $\alpha_i(t)$; color is held
fixed under deformation. $F_\theta$ first encodes $(\boldsymbol{\mu}_i, t)$
through a multi-resolution spatio-temporal feature grid (a HexPlane
comprising six learned 2D planes spanning each pairing of $\{x, y, z, t\}$,
bilinearly sampled and combined), and then decodes the resulting feature
through a shared multilayer perceptron trunk followed by four
output-specific heads, one per residual. This factorization allows the
network to share a common spatio-temporal representation across all four
deformed quantities while specializing each head's output.

\paragraph{Rendering.}
Given a camera pose $(\mathbf{R}, \mathbf{T})$ and time $t$, the deformed
Gaussians $\{(\boldsymbol{\mu}_i(t), \mathbf{s}_i(t), \mathbf{q}_i(t),
\alpha_i(t), \mathbf{c}_i)\}$ are projected to screen space and
alpha-composited in front-to-back order by a tile-based differentiable
rasterizer, producing a rendered color image $\hat{I}(\mathbf{R},
\mathbf{T}, t)$ and a per-pixel depth map $\hat{D}(\mathbf{R}, \mathbf{T},
t)$ at every pixel $\mathbf{p}$:
\begin{equation}
\hat{I}(\mathbf{p}) = \sum_{i \in \mathcal{N}(\mathbf{p})} \mathbf{c}_i\, \hat{\alpha}_i(\mathbf{p}) \prod_{j<i}\big(1-\hat{\alpha}_j(\mathbf{p})\big),
\end{equation}
where $\mathcal{N}(\mathbf{p})$ denotes the depth-ordered set of Gaussians
overlapping $\mathbf{p}$, and $\hat{\alpha}_i(\mathbf{p})$ is $\alpha_i(t)$
attenuated by the projected 2D footprint of $\boldsymbol{\Sigma}_i$ on the
image plane; $\hat{D}$ is defined as the analogous depth-weighted
composite.

\paragraph{Training.}
$\mathcal{G}$ and $\theta$ are optimized in two stages: a \emph{coarse}
stage that fits the canonical, undeformed Gaussians, followed by a
\emph{fine} stage that jointly optimizes $\mathcal{G}$ and $F_\theta$. The
base training objective combines a photometric term (an $\ell_1$ term and a
D-SSIM term against the ground-truth video frame), a depth term (a
scale-invariant inverse-depth $\ell_1$ loss or a Pearson-correlation loss
against the recorded depth, depending on the capture setup), a
total-variation term on the rendered image and depth that discourages
high-frequency noise, and, during the fine stage, a deformation regularizer
consisting of temporal-smoothness and total-variation penalties on the
HexPlane feature grid. Sections~\ref{sec:method:gcsr} and
\ref{sec:method:rtvc} extend and evaluate this base model, respectively.

\subsection{Geometry-Constrained Sideview Regularization (GCSR)}
\label{sec:method:gcsr}

Endoscopic video samples a narrow camera trajectory: the endoscope remains
close to a single working distance and rarely departs far from an
approximately frontal view of the tissue. Consequently, a Gaussian field
can be photometrically accurate along the recorded trajectory while
containing geometry that is implausible from any other viewpoint --
needle- or plate-shaped Gaussians and disconnected floaters that project to
a thin, inconsequential silhouette under the training views but yield a
discontinuous depth surface under even a small viewpoint perturbation.
Because free-viewpoint rendering is a design goal of this work, we
regularize the model against this failure mode during training, using
views synthesized by the model itself rather than any additional captured
data.

\paragraph{Sideview camera synthesis.}
At training iterations that are multiples of a fixed interval $\tau$, we
take the current training view's camera pose $(\mathbf{R}, \mathbf{T})$
and its recorded depth map, and compute the scene's working distance $d$
as the median of the valid, non-tool-masked depth values. An azimuth $a$
is then sampled, either uniformly from $[0^\circ, 360^\circ)$ or from a
fixed discrete set, and an elevation magnitude $e$ is fixed. A
\emph{sideview} camera pose $(\mathbf{R}_{a,e}, \mathbf{T}_{a,e})$ is
synthesized by rotating the original camera by angle $e$ about the axis
$\mathbf{n}(a) = (-\sin a, \cos a, 0)^\top$, pivoting about the point at
distance $d$ along the camera's principal ray. This defines an orbit of
the camera about the tissue surface at its own working distance, and thus
corresponds to a lateral perturbation of the endoscope rather than a
translation along the viewing axis. Fixing $\mathbf{n}(a) = (-\sin
0^\circ, \cos 0^\circ, 0)^\top$ and sweeping $e$ continuously through zero
into negative values reproduces the two opposing sideview directions
(nominally $a = 0^\circ$ and $a = 180^\circ$) as a single signed rotation
about one axis.

\paragraph{Depth-smoothness loss.}
The current Gaussians are rendered at the synthesized sideview pose to
obtain a depth map $\tilde{D}$; no ground truth is required at this pose.
Its local roughness is penalized with a Huber (smooth-$\ell_1$)
total-variation term:
\begin{equation}
\mathcal{L}_{\mathrm{depth}} = \sum_{(u,v)} h_\beta\!\big(\tilde{D}_{u+1,v} - \tilde{D}_{u,v}\big) + h_\beta\!\big(\tilde{D}_{u,v+1} - \tilde{D}_{u,v}\big),
\label{eq:gcsr-depth}
\end{equation}
\begin{equation}
h_\beta(x) = \begin{cases} \dfrac{x^2}{2\beta}, & |x| < \beta \\[4pt] |x| - \dfrac{\beta}{2}, & |x| \ge \beta. \end{cases}
\end{equation}
In contrast to a quadratic total-variation penalty, $h_\beta$ is quadratic,
and therefore strongly smoothing, for the sub-pixel depth noise that
dominates near $x = 0$, but grows only linearly beyond the threshold
$\beta$. This suppresses disconnected floaters that are visible only from
sideview poses without excessively penalizing genuine depth
discontinuities such as tissue folds. An analogous total-variation term on
the sideview color render, $\mathcal{L}_{\mathrm{smooth}}$, is available as
an option but is not the primary contribution of this work.

\paragraph{Anisotropy penalty.}
As a complement to the depth term above, we directly penalize the
needle- and plate-shaped Gaussians that are a common source of
view-dependent floaters, independent of viewpoint. For each Gaussian, an
anisotropy ratio $r_i = \max_k s_i^{(k)} / \big(\min_k s_i^{(k)} +
\epsilon\big)$ is defined and penalized as
\begin{equation}
p(r_i) = \begin{cases} (r_i - 1)^{\gamma}, & \text{no threshold} \\[2pt] \mathrm{relu}(r_i - r_0)^{\gamma}, & \text{threshold } r_0, \end{cases}
\end{equation}
\begin{equation}
\mathcal{L}_{\mathrm{aniso}} = \frac{1}{N}\sum_{i=1}^{N} p(r_i)\cdot \mathrm{sg}\!\big(\max_k s_i^{(k)}\big)^{\gamma_s},
\label{eq:gcsr-aniso}
\end{equation}
where $\mathrm{sg}(\cdot)$ denotes a stop-gradient operator. The
unthresholded form grows super-linearly with $r_i$ for $\gamma > 1$ and has
nonzero gradient everywhere, contracting the anisotropy distribution as a
whole; the thresholded form penalizes only ratios exceeding $r_0$, with
zero gradient below it. In both cases, the per-Gaussian weighting factor is
detached from the gradient and raised to a sub-linear power $\gamma_s < 1$,
so that the penalty reshapes anisotropic Gaussians without being dominated
by, or directly reducing the size of, the largest ones.

\paragraph{Total objective.}
The complete training loss augments the base objective of
Section~\ref{sec:method:3dgs} with both terms above, each active only when
its corresponding weight is nonzero:
\begin{equation}
\mathcal{L} = \mathcal{L}_{\mathrm{base}} + \lambda_{\mathrm{depth}}\,\mathcal{L}_{\mathrm{depth}} + \lambda_{\mathrm{smooth}}\,\mathcal{L}_{\mathrm{smooth}} + \lambda_{\mathrm{aniso}}\,\mathcal{L}_{\mathrm{aniso}}.
\end{equation}
The sideview terms are evaluated at a single, randomly sampled sideview
pose per training view every $\tau$ iterations rather than at every
iteration, limiting their computational overhead. No sideview camera is
synthesized at inference or rendering time, so GCSR introduces no
additional runtime cost.

\subsection{Round-Trip View Consistency (RTVC)}
\label{sec:method:rtvc}

Quantifying free-viewpoint reconstruction quality is inherently difficult
for monocular endoscopic video: computing PSNR or SSIM against a
photograph requires a captured image at the evaluation viewpoint, yet the
recorded video visits only the narrow surgical trajectory, so no such
image exists away from it. Round-Trip View Consistency (RTVC) addresses
this limitation without requiring any additional capture, by reusing the
ground truth already available at the original camera pose for the same
time instant, and evaluating whether the model's synthesized view at a
nearby off-axis pose is consistent with it through a round trip back
through that synthesized viewpoint.

\paragraph{Protocol.}
For a test frame with original pose $(\mathbf{R}, \mathbf{T})$ and ground
truth $I$, and for an evaluation elevation $e$ (and, optionally, one of the
four cardinal azimuths, left, right, up, or down), RTVC proceeds as
follows.
\begin{enumerate}
\item \textbf{Render.} The sideview pose $(\mathbf{R}_{a,e},
\mathbf{T}_{a,e})$ is synthesized as in Section~\ref{sec:method:gcsr}, and
the trained model is rendered at this pose to obtain a color image
$\tilde{I}$ and depth map $\tilde{D}$.
\item \textbf{Warp back.} $(\tilde{I}, \tilde{D})$ is forward-warped from
the sideview pose to the original pose using a depth-based, $z$-buffered
point splat: each valid sideview pixel is unprojected to a 3D point using
$\tilde{D}$ and the known sideview camera intrinsics and pose, reprojected
into the original camera, and rasterized with nearest-wins $z$-buffering to
resolve duplicate or occluded projections. This produces a warped image
$\hat{W}$ and a validity mask $M$ indicating which pixels of the original
frame the sideview render's geometry was able to re-explain.
\item \textbf{Score.} $\hat{W}$ is compared against the ground truth $I$
using masked PSNR and masked SSIM, restricted to pixels valid under both
$M$ and the surgical-tool mask, and, for SSIM, to windows whose valid
coverage exceeds a minimum threshold.
\end{enumerate}

\paragraph{Ground-truth-free property.}
No photograph at the sideview pose is required at any point in this
procedure: the only ground truth used is $I$, which was already captured
at the original pose. The resulting comparison is nonetheless a genuine
photometric evaluation, because any error in the reconstructed geometry at
the sideview pose, such as incorrect depth or missing or hallucinated
structure, displaces content relative to the true pixels of $I$ once
warped back, degrading RTVC's PSNR and SSIM even when the sideview render
$\tilde{I}$ appears plausible in isolation and cannot itself be validated
against a captured photograph.

\paragraph{Statistical evaluation.}
For a given elevation $e$, the up to four cardinal-azimuth measurements are
averaged per frame prior to any statistical comparison, as they share the
same underlying reconstruction at that time step and do not constitute
independent samples; treating them as independent would pseudo-replicate
the data and overstate statistical significance. This yields a single
paired sample per frame per model, which is compared between two models
using a paired $t$-test across frames. The special case $e = 0^\circ$,
corresponding to no sideview offset and no warp, reduces to the
conventional test-view PSNR and SSIM, and is reported alongside the RTVC
elevations as a reference point.

%
% Notation used throughout this section:
%   scalars      x
%   vectors      \mathbf{x}
%   matrices     \mathbf{X}
%   sets/scenes  \mathcal{X}

\section{Method}
\label{sec:method}

This section describes the proposed method in three parts. Section
\ref{sec:method:3dgs} presents the deformable 3D Gaussian Splatting
representation on which our approach is built. Section
\ref{sec:method:gcsr} introduces Geometry-Constrained Sideview
Regularization (GCSR), a training-time regularizer that constrains scene
geometry at viewpoints not present in the training trajectory. Section
\ref{sec:method:rtvc} introduces Round-Trip View Consistency (RTVC), an
evaluation protocol that quantifies reconstruction quality at such
viewpoints without requiring ground-truth imagery to be captured there.

\subsection{Deformable 3D Gaussian Splatting}
\label{sec:method:3dgs}

\paragraph{Canonical representation.}
The reconstructed tissue is represented as a set of $N$ canonical 3D
Gaussians $\mathcal{G} = \{(\boldsymbol{\mu}_i, \mathbf{s}_i, \mathbf{q}_i,
\alpha_i, \mathbf{c}_i)\}_{i=1}^{N}$, where $\boldsymbol{\mu}_i \in
\mathbb{R}^3$ denotes the mean, $\mathbf{s}_i \in \mathbb{R}^3_{>0}$ the
per-axis scale, $\mathbf{q}_i$ a unit quaternion specifying orientation,
$\alpha_i \in [0,1]$ an opacity, and $\mathbf{c}_i$ a set of
spherical-harmonics coefficients encoding view-dependent color. The
covariance of Gaussian $i$ is given by
\begin{equation}
\boldsymbol{\Sigma}_i = \mathbf{R}(\mathbf{q}_i)\,\mathrm{diag}(\mathbf{s}_i)^2\,\mathbf{R}(\mathbf{q}_i)^\top ,
\end{equation}
where $\mathbf{R}(\mathbf{q}_i)$ is the rotation matrix corresponding to
$\mathbf{q}_i$.

\paragraph{Deformation network.}
Because surgical tissue deforms over the course of a video, $\mathcal{G}$ is
not treated as static. A deformation network $F_\theta$ maps each canonical
Gaussian and a query time $t$ to a set of residuals,
\begin{equation}
\big(\Delta\boldsymbol{\mu}_i(t),\, \Delta\mathbf{s}_i(t),\, \Delta\mathbf{q}_i(t),\, \Delta\alpha_i(t)\big) = F_\theta(\boldsymbol{\mu}_i, t),
\end{equation}
which define the time-$t$ Gaussian $\boldsymbol{\mu}_i(t) =
\boldsymbol{\mu}_i + \Delta\boldsymbol{\mu}_i(t)$, and analogously for
$\mathbf{s}_i(t)$, $\mathbf{q}_i(t)$, and $\alpha_i(t)$; color is held
fixed under deformation. $F_\theta$ first encodes $(\boldsymbol{\mu}_i, t)$
through a multi-resolution spatio-temporal feature grid (a HexPlane
comprising six learned 2D planes spanning each pairing of $\{x, y, z, t\}$,
bilinearly sampled and combined), and then decodes the resulting feature
through a shared multilayer perceptron trunk followed by four
output-specific heads, one per residual. This factorization allows the
network to share a common spatio-temporal representation across all four
deformed quantities while specializing each head's output.

\paragraph{Rendering.}
Given a camera pose $(\mathbf{R}, \mathbf{T})$ and time $t$, the deformed
Gaussians $\{(\boldsymbol{\mu}_i(t), \mathbf{s}_i(t), \mathbf{q}_i(t),
\alpha_i(t), \mathbf{c}_i)\}$ are projected to screen space and
alpha-composited in front-to-back order by a tile-based differentiable
rasterizer, producing a rendered color image $\hat{I}(\mathbf{R},
\mathbf{T}, t)$ and a per-pixel depth map $\hat{D}(\mathbf{R}, \mathbf{T},
t)$ at every pixel $\mathbf{p}$:
\begin{equation}
\hat{I}(\mathbf{p}) = \sum_{i \in \mathcal{N}(\mathbf{p})} \mathbf{c}_i\, \hat{\alpha}_i(\mathbf{p}) \prod_{j<i}\big(1-\hat{\alpha}_j(\mathbf{p})\big),
\end{equation}
where $\mathcal{N}(\mathbf{p})$ denotes the depth-ordered set of Gaussians
overlapping $\mathbf{p}$, and $\hat{\alpha}_i(\mathbf{p})$ is $\alpha_i(t)$
attenuated by the projected 2D footprint of $\boldsymbol{\Sigma}_i$ on the
image plane; $\hat{D}$ is defined as the analogous depth-weighted
composite.

\paragraph{Training.}
$\mathcal{G}$ and $\theta$ are optimized in two stages: a \emph{coarse}
stage that fits the canonical, undeformed Gaussians, followed by a
\emph{fine} stage that jointly optimizes $\mathcal{G}$ and $F_\theta$. The
base training objective combines a photometric term (an $\ell_1$ term and a
D-SSIM term against the ground-truth video frame), a depth term (a
scale-invariant inverse-depth $\ell_1$ loss or a Pearson-correlation loss
against the recorded depth, depending on the capture setup), a
total-variation term on the rendered image and depth that discourages
high-frequency noise, and, during the fine stage, a deformation regularizer
consisting of temporal-smoothness and total-variation penalties on the
HexPlane feature grid. Sections~\ref{sec:method:gcsr} and
\ref{sec:method:rtvc} extend and evaluate this base model, respectively.

\subsection{Geometry-Constrained Sideview Regularization (GCSR)}
\label{sec:method:gcsr}

Endoscopic video samples a narrow camera trajectory: the endoscope remains
close to a single working distance and rarely departs far from an
approximately frontal view of the tissue. Consequently, a Gaussian field
can be photometrically accurate along the recorded trajectory while
containing geometry that is implausible from any other viewpoint --
needle- or plate-shaped Gaussians and disconnected floaters that project to
a thin, inconsequential silhouette under the training views but yield a
discontinuous depth surface under even a small viewpoint perturbation.
Because free-viewpoint rendering is a design goal of this work, we
regularize the model against this failure mode during training, using
views synthesized by the model itself rather than any additional captured
data.

\paragraph{Sideview camera synthesis.}
At training iterations that are multiples of a fixed interval $\tau$, we
take the current training view's camera pose $(\mathbf{R}, \mathbf{T})$
and its recorded depth map, and compute the scene's working distance $d$
as the median of the valid, non-tool-masked depth values. An azimuth $a$
is then sampled, either uniformly from $[0^\circ, 360^\circ)$ or from a
fixed discrete set, and an elevation magnitude $e$ is fixed. A
\emph{sideview} camera pose $(\mathbf{R}_{a,e}, \mathbf{T}_{a,e})$ is
synthesized by rotating the original camera by angle $e$ about the axis
$\mathbf{n}(a) = (-\sin a, \cos a, 0)^\top$, pivoting about the point at
distance $d$ along the camera's principal ray. This defines an orbit of
the camera about the tissue surface at its own working distance, and thus
corresponds to a lateral perturbation of the endoscope rather than a
translation along the viewing axis. Fixing $\mathbf{n}(a) = (-\sin
0^\circ, \cos 0^\circ, 0)^\top$ and sweeping $e$ continuously through zero
into negative values reproduces the two opposing sideview directions
(nominally $a = 0^\circ$ and $a = 180^\circ$) as a single signed rotation
about one axis.

\paragraph{Depth-smoothness loss.}
The current Gaussians are rendered at the synthesized sideview pose to
obtain a depth map $\tilde{D}$; no ground truth is required at this pose.
Its local roughness is penalized with a Huber (smooth-$\ell_1$)
total-variation term:
\begin{equation}
\mathcal{L}_{\mathrm{depth}} = \sum_{(u,v)} h_\beta\!\big(\tilde{D}_{u+1,v} - \tilde{D}_{u,v}\big) + h_\beta\!\big(\tilde{D}_{u,v+1} - \tilde{D}_{u,v}\big),
\label{eq:gcsr-depth}
\end{equation}
\begin{equation}
h_\beta(x) = \begin{cases} \dfrac{x^2}{2\beta}, & |x| < \beta \\[4pt] |x| - \dfrac{\beta}{2}, & |x| \ge \beta. \end{cases}
\end{equation}
In contrast to a quadratic total-variation penalty, $h_\beta$ is quadratic,
and therefore strongly smoothing, for the sub-pixel depth noise that
dominates near $x = 0$, but grows only linearly beyond the threshold
$\beta$. This suppresses disconnected floaters that are visible only from
sideview poses without excessively penalizing genuine depth
discontinuities such as tissue folds. An analogous total-variation term on
the sideview color render, $\mathcal{L}_{\mathrm{smooth}}$, is available as
an option but is not the primary contribution of this work.

\paragraph{Anisotropy penalty.}
As a complement to the depth term above, we directly penalize the
needle- and plate-shaped Gaussians that are a common source of
view-dependent floaters, independent of viewpoint. For each Gaussian, an
anisotropy ratio $r_i = \max_k s_i^{(k)} / \big(\min_k s_i^{(k)} +
\epsilon\big)$ is defined and penalized as
\begin{equation}
p(r_i) = \begin{cases} (r_i - 1)^{\gamma}, & \text{no threshold} \\[2pt] \mathrm{relu}(r_i - r_0)^{\gamma}, & \text{threshold } r_0, \end{cases}
\end{equation}
\begin{equation}
\mathcal{L}_{\mathrm{aniso}} = \frac{1}{N}\sum_{i=1}^{N} p(r_i)\cdot \mathrm{sg}\!\big(\max_k s_i^{(k)}\big)^{\gamma_s},
\label{eq:gcsr-aniso}
\end{equation}
where $\mathrm{sg}(\cdot)$ denotes a stop-gradient operator. The
unthresholded form grows super-linearly with $r_i$ for $\gamma > 1$ and has
nonzero gradient everywhere, contracting the anisotropy distribution as a
whole; the thresholded form penalizes only ratios exceeding $r_0$, with
zero gradient below it. In both cases, the per-Gaussian weighting factor is
detached from the gradient and raised to a sub-linear power $\gamma_s < 1$,
so that the penalty reshapes anisotropic Gaussians without being dominated
by, or directly reducing the size of, the largest ones.

\paragraph{Total objective.}
The complete training loss augments the base objective of
Section~\ref{sec:method:3dgs} with both terms above, each active only when
its corresponding weight is nonzero:
\begin{equation}
\mathcal{L} = \mathcal{L}_{\mathrm{base}} + \lambda_{\mathrm{depth}}\,\mathcal{L}_{\mathrm{depth}} + \lambda_{\mathrm{smooth}}\,\mathcal{L}_{\mathrm{smooth}} + \lambda_{\mathrm{aniso}}\,\mathcal{L}_{\mathrm{aniso}}.
\end{equation}
The sideview terms are evaluated at a single, randomly sampled sideview
pose per training view every $\tau$ iterations rather than at every
iteration, limiting their computational overhead. No sideview camera is
synthesized at inference or rendering time, so GCSR introduces no
additional runtime cost.

\subsection{Round-Trip View Consistency (RTVC)}
\label{sec:method:rtvc}

Quantifying free-viewpoint reconstruction quality is inherently difficult
for monocular endoscopic video: computing PSNR or SSIM against a
photograph requires a captured image at the evaluation viewpoint, yet the
recorded video visits only the narrow surgical trajectory, so no such
image exists away from it. Round-Trip View Consistency (RTVC) addresses
this limitation without requiring any additional capture, by reusing the
ground truth already available at the original camera pose for the same
time instant, and evaluating whether the model's synthesized view at a
nearby off-axis pose is consistent with it through a round trip back
through that synthesized viewpoint.

\paragraph{Protocol.}
For a test frame with original pose $(\mathbf{R}, \mathbf{T})$ and ground
truth $I$, and for an evaluation elevation $e$ (and, optionally, one of the
four cardinal azimuths, left, right, up, or down), RTVC proceeds as
follows.
\begin{enumerate}
\item \textbf{Render.} The sideview pose $(\mathbf{R}_{a,e},
\mathbf{T}_{a,e})$ is synthesized as in Section~\ref{sec:method:gcsr}, and
the trained model is rendered at this pose to obtain a color image
$\tilde{I}$ and depth map $\tilde{D}$.
\item \textbf{Warp back.} $(\tilde{I}, \tilde{D})$ is forward-warped from
the sideview pose to the original pose using a depth-based, $z$-buffered
point splat: each valid sideview pixel is unprojected to a 3D point using
$\tilde{D}$ and the known sideview camera intrinsics and pose, reprojected
into the original camera, and rasterized with nearest-wins $z$-buffering to
resolve duplicate or occluded projections. This produces a warped image
$\hat{W}$ and a validity mask $M$ indicating which pixels of the original
frame the sideview render's geometry was able to re-explain.
\item \textbf{Score.} $\hat{W}$ is compared against the ground truth $I$
using masked PSNR and masked SSIM, restricted to pixels valid under both
$M$ and the surgical-tool mask, and, for SSIM, to windows whose valid
coverage exceeds a minimum threshold.
\end{enumerate}

\paragraph{Ground-truth-free property.}
No photograph at the sideview pose is required at any point in this
procedure: the only ground truth used is $I$, which was already captured
at the original pose. The resulting comparison is nonetheless a genuine
photometric evaluation, because any error in the reconstructed geometry at
the sideview pose, such as incorrect depth or missing or hallucinated
structure, displaces content relative to the true pixels of $I$ once
warped back, degrading RTVC's PSNR and SSIM even when the sideview render
$\tilde{I}$ appears plausible in isolation and cannot itself be validated
against a captured photograph.

\paragraph{Statistical evaluation.}
For a given elevation $e$, the up to four cardinal-azimuth measurements are
averaged per frame prior to any statistical comparison, as they share the
same underlying reconstruction at that time step and do not constitute
independent samples; treating them as independent would pseudo-replicate
the data and overstate statistical significance. This yields a single
paired sample per frame per model, which is compared between two models
using a paired $t$-test across frames. The special case $e = 0^\circ$,
corresponding to no sideview offset and no warp, reduces to the
conventional test-view PSNR and SSIM, and is reported alongside the RTVC
elevations as a reference point.

%
% Notation used throughout this section:
%   scalars      x
%   vectors      \mathbf{x}
%   matrices     \mathbf{X}
%   sets/scenes  \mathcal{X}

\section{Method}
\label{sec:method}

This section describes the proposed method in three parts. Section
\ref{sec:method:3dgs} presents the deformable 3D Gaussian Splatting
representation on which our approach is built. Section
\ref{sec:method:gcsr} introduces Geometry-Constrained Sideview
Regularization (GCSR), a training-time regularizer that constrains scene
geometry at viewpoints not present in the training trajectory. Section
\ref{sec:method:rtvc} introduces Round-Trip View Consistency (RTVC), an
evaluation protocol that quantifies reconstruction quality at such
viewpoints without requiring ground-truth imagery to be captured there.

\subsection{Deformable 3D Gaussian Splatting}
\label{sec:method:3dgs}

\paragraph{Canonical representation.}
The reconstructed tissue is represented as a set of $N$ canonical 3D
Gaussians $\mathcal{G} = \{(\boldsymbol{\mu}_i, \mathbf{s}_i, \mathbf{q}_i,
\alpha_i, \mathbf{c}_i)\}_{i=1}^{N}$, where $\boldsymbol{\mu}_i \in
\mathbb{R}^3$ denotes the mean, $\mathbf{s}_i \in \mathbb{R}^3_{>0}$ the
per-axis scale, $\mathbf{q}_i$ a unit quaternion specifying orientation,
$\alpha_i \in [0,1]$ an opacity, and $\mathbf{c}_i$ a set of
spherical-harmonics coefficients encoding view-dependent color. The
covariance of Gaussian $i$ is given by
\begin{equation}
\boldsymbol{\Sigma}_i = \mathbf{R}(\mathbf{q}_i)\,\mathrm{diag}(\mathbf{s}_i)^2\,\mathbf{R}(\mathbf{q}_i)^\top ,
\end{equation}
where $\mathbf{R}(\mathbf{q}_i)$ is the rotation matrix corresponding to
$\mathbf{q}_i$.

\paragraph{Deformation network.}
Because surgical tissue deforms over the course of a video, $\mathcal{G}$ is
not treated as static. A deformation network $F_\theta$ maps each canonical
Gaussian and a query time $t$ to a set of residuals,
\begin{equation}
\big(\Delta\boldsymbol{\mu}_i(t),\, \Delta\mathbf{s}_i(t),\, \Delta\mathbf{q}_i(t),\, \Delta\alpha_i(t)\big) = F_\theta(\boldsymbol{\mu}_i, t),
\end{equation}
which define the time-$t$ Gaussian $\boldsymbol{\mu}_i(t) =
\boldsymbol{\mu}_i + \Delta\boldsymbol{\mu}_i(t)$, and analogously for
$\mathbf{s}_i(t)$, $\mathbf{q}_i(t)$, and $\alpha_i(t)$; color is held
fixed under deformation. $F_\theta$ first encodes $(\boldsymbol{\mu}_i, t)$
through a multi-resolution spatio-temporal feature grid (a HexPlane
comprising six learned 2D planes spanning each pairing of $\{x, y, z, t\}$,
bilinearly sampled and combined), and then decodes the resulting feature
through a shared multilayer perceptron trunk followed by four
output-specific heads, one per residual. This factorization allows the
network to share a common spatio-temporal representation across all four
deformed quantities while specializing each head's output.

\paragraph{Rendering.}
Given a camera pose $(\mathbf{R}, \mathbf{T})$ and time $t$, the deformed
Gaussians $\{(\boldsymbol{\mu}_i(t), \mathbf{s}_i(t), \mathbf{q}_i(t),
\alpha_i(t), \mathbf{c}_i)\}$ are projected to screen space and
alpha-composited in front-to-back order by a tile-based differentiable
rasterizer, producing a rendered color image $\hat{I}(\mathbf{R},
\mathbf{T}, t)$ and a per-pixel depth map $\hat{D}(\mathbf{R}, \mathbf{T},
t)$ at every pixel $\mathbf{p}$:
\begin{equation}
\hat{I}(\mathbf{p}) = \sum_{i \in \mathcal{N}(\mathbf{p})} \mathbf{c}_i\, \hat{\alpha}_i(\mathbf{p}) \prod_{j<i}\big(1-\hat{\alpha}_j(\mathbf{p})\big),
\end{equation}
where $\mathcal{N}(\mathbf{p})$ denotes the depth-ordered set of Gaussians
overlapping $\mathbf{p}$, and $\hat{\alpha}_i(\mathbf{p})$ is $\alpha_i(t)$
attenuated by the projected 2D footprint of $\boldsymbol{\Sigma}_i$ on the
image plane; $\hat{D}$ is defined as the analogous depth-weighted
composite.

\paragraph{Training.}
$\mathcal{G}$ and $\theta$ are optimized in two stages: a \emph{coarse}
stage that fits the canonical, undeformed Gaussians, followed by a
\emph{fine} stage that jointly optimizes $\mathcal{G}$ and $F_\theta$. The
base training objective combines a photometric term (an $\ell_1$ term and a
D-SSIM term against the ground-truth video frame), a depth term (a
scale-invariant inverse-depth $\ell_1$ loss or a Pearson-correlation loss
against the recorded depth, depending on the capture setup), a
total-variation term on the rendered image and depth that discourages
high-frequency noise, and, during the fine stage, a deformation regularizer
consisting of temporal-smoothness and total-variation penalties on the
HexPlane feature grid. Sections~\ref{sec:method:gcsr} and
\ref{sec:method:rtvc} extend and evaluate this base model, respectively.

\subsection{Geometry-Constrained Sideview Regularization (GCSR)}
\label{sec:method:gcsr}

Endoscopic video samples a narrow camera trajectory: the endoscope remains
close to a single working distance and rarely departs far from an
approximately frontal view of the tissue. Consequently, a Gaussian field
can be photometrically accurate along the recorded trajectory while
containing geometry that is implausible from any other viewpoint --
needle- or plate-shaped Gaussians and disconnected floaters that project to
a thin, inconsequential silhouette under the training views but yield a
discontinuous depth surface under even a small viewpoint perturbation.
Because free-viewpoint rendering is a design goal of this work, we
regularize the model against this failure mode during training, using
views synthesized by the model itself rather than any additional captured
data.

\paragraph{Sideview camera synthesis.}
At training iterations that are multiples of a fixed interval $\tau$, we
take the current training view's camera pose $(\mathbf{R}, \mathbf{T})$
and its recorded depth map, and compute the scene's working distance $d$
as the median of the valid, non-tool-masked depth values. An azimuth $a$
is then sampled, either uniformly from $[0^\circ, 360^\circ)$ or from a
fixed discrete set, and an elevation magnitude $e$ is fixed. A
\emph{sideview} camera pose $(\mathbf{R}_{a,e}, \mathbf{T}_{a,e})$ is
synthesized by rotating the original camera by angle $e$ about the axis
$\mathbf{n}(a) = (-\sin a, \cos a, 0)^\top$, pivoting about the point at
distance $d$ along the camera's principal ray. This defines an orbit of
the camera about the tissue surface at its own working distance, and thus
corresponds to a lateral perturbation of the endoscope rather than a
translation along the viewing axis. Fixing $\mathbf{n}(a) = (-\sin
0^\circ, \cos 0^\circ, 0)^\top$ and sweeping $e$ continuously through zero
into negative values reproduces the two opposing sideview directions
(nominally $a = 0^\circ$ and $a = 180^\circ$) as a single signed rotation
about one axis.

\paragraph{Depth-smoothness loss.}
The current Gaussians are rendered at the synthesized sideview pose to
obtain a depth map $\tilde{D}$; no ground truth is required at this pose.
Its local roughness is penalized with a Huber (smooth-$\ell_1$)
total-variation term:
\begin{equation}
\mathcal{L}_{\mathrm{depth}} = \sum_{(u,v)} h_\beta\!\big(\tilde{D}_{u+1,v} - \tilde{D}_{u,v}\big) + h_\beta\!\big(\tilde{D}_{u,v+1} - \tilde{D}_{u,v}\big),
\label{eq:gcsr-depth}
\end{equation}
\begin{equation}
h_\beta(x) = \begin{cases} \dfrac{x^2}{2\beta}, & |x| < \beta \\[4pt] |x| - \dfrac{\beta}{2}, & |x| \ge \beta. \end{cases}
\end{equation}
In contrast to a quadratic total-variation penalty, $h_\beta$ is quadratic,
and therefore strongly smoothing, for the sub-pixel depth noise that
dominates near $x = 0$, but grows only linearly beyond the threshold
$\beta$. This suppresses disconnected floaters that are visible only from
sideview poses without excessively penalizing genuine depth
discontinuities such as tissue folds. An analogous total-variation term on
the sideview color render, $\mathcal{L}_{\mathrm{smooth}}$, is available as
an option but is not the primary contribution of this work.

\paragraph{Anisotropy penalty.}
As a complement to the depth term above, we directly penalize the
needle- and plate-shaped Gaussians that are a common source of
view-dependent floaters, independent of viewpoint. For each Gaussian, an
anisotropy ratio $r_i = \max_k s_i^{(k)} / \big(\min_k s_i^{(k)} +
\epsilon\big)$ is defined and penalized as
\begin{equation}
p(r_i) = \begin{cases} (r_i - 1)^{\gamma}, & \text{no threshold} \\[2pt] \mathrm{relu}(r_i - r_0)^{\gamma}, & \text{threshold } r_0, \end{cases}
\end{equation}
\begin{equation}
\mathcal{L}_{\mathrm{aniso}} = \frac{1}{N}\sum_{i=1}^{N} p(r_i)\cdot \mathrm{sg}\!\big(\max_k s_i^{(k)}\big)^{\gamma_s},
\label{eq:gcsr-aniso}
\end{equation}
where $\mathrm{sg}(\cdot)$ denotes a stop-gradient operator. The
unthresholded form grows super-linearly with $r_i$ for $\gamma > 1$ and has
nonzero gradient everywhere, contracting the anisotropy distribution as a
whole; the thresholded form penalizes only ratios exceeding $r_0$, with
zero gradient below it. In both cases, the per-Gaussian weighting factor is
detached from the gradient and raised to a sub-linear power $\gamma_s < 1$,
so that the penalty reshapes anisotropic Gaussians without being dominated
by, or directly reducing the size of, the largest ones.

\paragraph{Total objective.}
The complete training loss augments the base objective of
Section~\ref{sec:method:3dgs} with both terms above, each active only when
its corresponding weight is nonzero:
\begin{equation}
\mathcal{L} = \mathcal{L}_{\mathrm{base}} + \lambda_{\mathrm{depth}}\,\mathcal{L}_{\mathrm{depth}} + \lambda_{\mathrm{smooth}}\,\mathcal{L}_{\mathrm{smooth}} + \lambda_{\mathrm{aniso}}\,\mathcal{L}_{\mathrm{aniso}}.
\end{equation}
The sideview terms are evaluated at a single, randomly sampled sideview
pose per training view every $\tau$ iterations rather than at every
iteration, limiting their computational overhead. No sideview camera is
synthesized at inference or rendering time, so GCSR introduces no
additional runtime cost.

\subsection{Round-Trip View Consistency (RTVC)}
\label{sec:method:rtvc}

Quantifying free-viewpoint reconstruction quality is inherently difficult
for monocular endoscopic video: computing PSNR or SSIM against a
photograph requires a captured image at the evaluation viewpoint, yet the
recorded video visits only the narrow surgical trajectory, so no such
image exists away from it. Round-Trip View Consistency (RTVC) addresses
this limitation without requiring any additional capture, by reusing the
ground truth already available at the original camera pose for the same
time instant, and evaluating whether the model's synthesized view at a
nearby off-axis pose is consistent with it through a round trip back
through that synthesized viewpoint.

\paragraph{Protocol.}
For a test frame with original pose $(\mathbf{R}, \mathbf{T})$ and ground
truth $I$, and for an evaluation elevation $e$ (and, optionally, one of the
four cardinal azimuths, left, right, up, or down), RTVC proceeds as
follows.
\begin{enumerate}
\item \textbf{Render.} The sideview pose $(\mathbf{R}_{a,e},
\mathbf{T}_{a,e})$ is synthesized as in Section~\ref{sec:method:gcsr}, and
the trained model is rendered at this pose to obtain a color image
$\tilde{I}$ and depth map $\tilde{D}$.
\item \textbf{Warp back.} $(\tilde{I}, \tilde{D})$ is forward-warped from
the sideview pose to the original pose using a depth-based, $z$-buffered
point splat: each valid sideview pixel is unprojected to a 3D point using
$\tilde{D}$ and the known sideview camera intrinsics and pose, reprojected
into the original camera, and rasterized with nearest-wins $z$-buffering to
resolve duplicate or occluded projections. This produces a warped image
$\hat{W}$ and a validity mask $M$ indicating which pixels of the original
frame the sideview render's geometry was able to re-explain.
\item \textbf{Score.} $\hat{W}$ is compared against the ground truth $I$
using masked PSNR and masked SSIM, restricted to pixels valid under both
$M$ and the surgical-tool mask, and, for SSIM, to windows whose valid
coverage exceeds a minimum threshold.
\end{enumerate}

\paragraph{Ground-truth-free property.}
No photograph at the sideview pose is required at any point in this
procedure: the only ground truth used is $I$, which was already captured
at the original pose. The resulting comparison is nonetheless a genuine
photometric evaluation, because any error in the reconstructed geometry at
the sideview pose, such as incorrect depth or missing or hallucinated
structure, displaces content relative to the true pixels of $I$ once
warped back, degrading RTVC's PSNR and SSIM even when the sideview render
$\tilde{I}$ appears plausible in isolation and cannot itself be validated
against a captured photograph.

\paragraph{Statistical evaluation.}
For a given elevation $e$, the up to four cardinal-azimuth measurements are
averaged per frame prior to any statistical comparison, as they share the
same underlying reconstruction at that time step and do not constitute
independent samples; treating them as independent would pseudo-replicate
the data and overstate statistical significance. This yields a single
paired sample per frame per model, which is compared between two models
using a paired $t$-test across frames. The special case $e = 0^\circ$,
corresponding to no sideview offset and no warp, reduces to the
conventional test-view PSNR and SSIM, and is reported alongside the RTVC
elevations as a reference point.

%
% Notation used throughout this section:
%   scalars      x
%   vectors      \mathbf{x}
%   matrices     \mathbf{X}
%   sets/scenes  \mathcal{X}

\section{Method}
\label{sec:method}

This section describes the proposed method in three parts. Section
\ref{sec:method:3dgs} presents the deformable 3D Gaussian Splatting
representation on which our approach is built. Section
\ref{sec:method:gcsr} introduces Geometry-Constrained Sideview
Regularization (GCSR), a training-time regularizer that constrains scene
geometry at viewpoints not present in the training trajectory. Section
\ref{sec:method:rtvc} introduces Round-Trip View Consistency (RTVC), an
evaluation protocol that quantifies reconstruction quality at such
viewpoints without requiring ground-truth imagery to be captured there.

\subsection{Deformable 3D Gaussian Splatting}
\label{sec:method:3dgs}

\paragraph{Canonical representation.}
The reconstructed tissue is represented as a set of $N$ canonical 3D
Gaussians $\mathcal{G} = \{(\boldsymbol{\mu}_i, \mathbf{s}_i, \mathbf{q}_i,
\alpha_i, \mathbf{c}_i)\}_{i=1}^{N}$, where $\boldsymbol{\mu}_i \in
\mathbb{R}^3$ denotes the mean, $\mathbf{s}_i \in \mathbb{R}^3_{>0}$ the
per-axis scale, $\mathbf{q}_i$ a unit quaternion specifying orientation,
$\alpha_i \in [0,1]$ an opacity, and $\mathbf{c}_i$ a set of
spherical-harmonics coefficients encoding view-dependent color. The
covariance of Gaussian $i$ is given by
\begin{equation}
\boldsymbol{\Sigma}_i = \mathbf{R}(\mathbf{q}_i)\,\mathrm{diag}(\mathbf{s}_i)^2\,\mathbf{R}(\mathbf{q}_i)^\top ,
\end{equation}
where $\mathbf{R}(\mathbf{q}_i)$ is the rotation matrix corresponding to
$\mathbf{q}_i$.

\paragraph{Deformation network.}
Because surgical tissue deforms over the course of a video, $\mathcal{G}$ is
not treated as static. A deformation network $F_\theta$ maps each canonical
Gaussian and a query time $t$ to a set of residuals,
\begin{equation}
\big(\Delta\boldsymbol{\mu}_i(t),\, \Delta\mathbf{s}_i(t),\, \Delta\mathbf{q}_i(t),\, \Delta\alpha_i(t)\big) = F_\theta(\boldsymbol{\mu}_i, t),
\end{equation}
which define the time-$t$ Gaussian $\boldsymbol{\mu}_i(t) =
\boldsymbol{\mu}_i + \Delta\boldsymbol{\mu}_i(t)$, and analogously for
$\mathbf{s}_i(t)$, $\mathbf{q}_i(t)$, and $\alpha_i(t)$; color is held
fixed under deformation. $F_\theta$ first encodes $(\boldsymbol{\mu}_i, t)$
through a multi-resolution spatio-temporal feature grid (a HexPlane
comprising six learned 2D planes spanning each pairing of $\{x, y, z, t\}$,
bilinearly sampled and combined), and then decodes the resulting feature
through a shared multilayer perceptron trunk followed by four
output-specific heads, one per residual. This factorization allows the
network to share a common spatio-temporal representation across all four
deformed quantities while specializing each head's output.

\paragraph{Rendering.}
Given a camera pose $(\mathbf{R}, \mathbf{T})$ and time $t$, the deformed
Gaussians $\{(\boldsymbol{\mu}_i(t), \mathbf{s}_i(t), \mathbf{q}_i(t),
\alpha_i(t), \mathbf{c}_i)\}$ are projected to screen space and
alpha-composited in front-to-back order by a tile-based differentiable
rasterizer, producing a rendered color image $\hat{I}(\mathbf{R},
\mathbf{T}, t)$ and a per-pixel depth map $\hat{D}(\mathbf{R}, \mathbf{T},
t)$ at every pixel $\mathbf{p}$:
\begin{equation}
\hat{I}(\mathbf{p}) = \sum_{i \in \mathcal{N}(\mathbf{p})} \mathbf{c}_i\, \hat{\alpha}_i(\mathbf{p}) \prod_{j<i}\big(1-\hat{\alpha}_j(\mathbf{p})\big),
\end{equation}
where $\mathcal{N}(\mathbf{p})$ denotes the depth-ordered set of Gaussians
overlapping $\mathbf{p}$, and $\hat{\alpha}_i(\mathbf{p})$ is $\alpha_i(t)$
attenuated by the projected 2D footprint of $\boldsymbol{\Sigma}_i$ on the
image plane; $\hat{D}$ is defined as the analogous depth-weighted
composite.

\paragraph{Training.}
$\mathcal{G}$ and $\theta$ are optimized in two stages: a \emph{coarse}
stage that fits the canonical, undeformed Gaussians, followed by a
\emph{fine} stage that jointly optimizes $\mathcal{G}$ and $F_\theta$. The
base training objective combines a photometric term (an $\ell_1$ term and a
D-SSIM term against the ground-truth video frame), a depth term (a
scale-invariant inverse-depth $\ell_1$ loss or a Pearson-correlation loss
against the recorded depth, depending on the capture setup), a
total-variation term on the rendered image and depth that discourages
high-frequency noise, and, during the fine stage, a deformation regularizer
consisting of temporal-smoothness and total-variation penalties on the
HexPlane feature grid. Sections~\ref{sec:method:gcsr} and
\ref{sec:method:rtvc} extend and evaluate this base model, respectively.

\subsection{Geometry-Constrained Sideview Regularization (GCSR)}
\label{sec:method:gcsr}

Endoscopic video samples a narrow camera trajectory: the endoscope remains
close to a single working distance and rarely departs far from an
approximately frontal view of the tissue. Consequently, a Gaussian field
can be photometrically accurate along the recorded trajectory while
containing geometry that is implausible from any other viewpoint --
needle- or plate-shaped Gaussians and disconnected floaters that project to
a thin, inconsequential silhouette under the training views but yield a
discontinuous depth surface under even a small viewpoint perturbation.
Because free-viewpoint rendering is a design goal of this work, we
regularize the model against this failure mode during training, using
views synthesized by the model itself rather than any additional captured
data.

\paragraph{Sideview camera synthesis.}
At training iterations that are multiples of a fixed interval $\tau$, we
take the current training view's camera pose $(\mathbf{R}, \mathbf{T})$
and its recorded depth map, and compute the scene's working distance $d$
as the median of the valid, non-tool-masked depth values. An azimuth $a$
is then sampled, either uniformly from $[0^\circ, 360^\circ)$ or from a
fixed discrete set, and an elevation magnitude $e$ is fixed. A
\emph{sideview} camera pose $(\mathbf{R}_{a,e}, \mathbf{T}_{a,e})$ is
synthesized by rotating the original camera by angle $e$ about the axis
$\mathbf{n}(a) = (-\sin a, \cos a, 0)^\top$, pivoting about the point at
distance $d$ along the camera's principal ray. This defines an orbit of
the camera about the tissue surface at its own working distance, and thus
corresponds to a lateral perturbation of the endoscope rather than a
translation along the viewing axis. Fixing $\mathbf{n}(a) = (-\sin
0^\circ, \cos 0^\circ, 0)^\top$ and sweeping $e$ continuously through zero
into negative values reproduces the two opposing sideview directions
(nominally $a = 0^\circ$ and $a = 180^\circ$) as a single signed rotation
about one axis.

\paragraph{Depth-smoothness loss.}
The current Gaussians are rendered at the synthesized sideview pose to
obtain a depth map $\tilde{D}$; no ground truth is required at this pose.
Its local roughness is penalized with a Huber (smooth-$\ell_1$)
total-variation term:
\begin{equation}
\mathcal{L}_{\mathrm{depth}} = \sum_{(u,v)} h_\beta\!\big(\tilde{D}_{u+1,v} - \tilde{D}_{u,v}\big) + h_\beta\!\big(\tilde{D}_{u,v+1} - \tilde{D}_{u,v}\big),
\label{eq:gcsr-depth}
\end{equation}
\begin{equation}
h_\beta(x) = \begin{cases} \dfrac{x^2}{2\beta}, & |x| < \beta \\[4pt] |x| - \dfrac{\beta}{2}, & |x| \ge \beta. \end{cases}
\end{equation}
In contrast to a quadratic total-variation penalty, $h_\beta$ is quadratic,
and therefore strongly smoothing, for the sub-pixel depth noise that
dominates near $x = 0$, but grows only linearly beyond the threshold
$\beta$. This suppresses disconnected floaters that are visible only from
sideview poses without excessively penalizing genuine depth
discontinuities such as tissue folds. An analogous total-variation term on
the sideview color render, $\mathcal{L}_{\mathrm{smooth}}$, is available as
an option but is not the primary contribution of this work.

\paragraph{Anisotropy penalty.}
As a complement to the depth term above, we directly penalize the
needle- and plate-shaped Gaussians that are a common source of
view-dependent floaters, independent of viewpoint. For each Gaussian, an
anisotropy ratio $r_i = \max_k s_i^{(k)} / \big(\min_k s_i^{(k)} +
\epsilon\big)$ is defined and penalized as
\begin{equation}
p(r_i) = \begin{cases} (r_i - 1)^{\gamma}, & \text{no threshold} \\[2pt] \mathrm{relu}(r_i - r_0)^{\gamma}, & \text{threshold } r_0, \end{cases}
\end{equation}
\begin{equation}
\mathcal{L}_{\mathrm{aniso}} = \frac{1}{N}\sum_{i=1}^{N} p(r_i)\cdot \mathrm{sg}\!\big(\max_k s_i^{(k)}\big)^{\gamma_s},
\label{eq:gcsr-aniso}
\end{equation}
where $\mathrm{sg}(\cdot)$ denotes a stop-gradient operator. The
unthresholded form grows super-linearly with $r_i$ for $\gamma > 1$ and has
nonzero gradient everywhere, contracting the anisotropy distribution as a
whole; the thresholded form penalizes only ratios exceeding $r_0$, with
zero gradient below it. In both cases, the per-Gaussian weighting factor is
detached from the gradient and raised to a sub-linear power $\gamma_s < 1$,
so that the penalty reshapes anisotropic Gaussians without being dominated
by, or directly reducing the size of, the largest ones.

\paragraph{Total objective.}
The complete training loss augments the base objective of
Section~\ref{sec:method:3dgs} with both terms above, each active only when
its corresponding weight is nonzero:
\begin{equation}
\mathcal{L} = \mathcal{L}_{\mathrm{base}} + \lambda_{\mathrm{depth}}\,\mathcal{L}_{\mathrm{depth}} + \lambda_{\mathrm{smooth}}\,\mathcal{L}_{\mathrm{smooth}} + \lambda_{\mathrm{aniso}}\,\mathcal{L}_{\mathrm{aniso}}.
\end{equation}
The sideview terms are evaluated at a single, randomly sampled sideview
pose per training view every $\tau$ iterations rather than at every
iteration, limiting their computational overhead. No sideview camera is
synthesized at inference or rendering time, so GCSR introduces no
additional runtime cost.

\subsection{Round-Trip View Consistency (RTVC)}
\label{sec:method:rtvc}

Quantifying free-viewpoint reconstruction quality is inherently difficult
for monocular endoscopic video: computing PSNR or SSIM against a
photograph requires a captured image at the evaluation viewpoint, yet the
recorded video visits only the narrow surgical trajectory, so no such
image exists away from it. Round-Trip View Consistency (RTVC) addresses
this limitation without requiring any additional capture, by reusing the
ground truth already available at the original camera pose for the same
time instant, and evaluating whether the model's synthesized view at a
nearby off-axis pose is consistent with it through a round trip back
through that synthesized viewpoint.

\paragraph{Protocol.}
For a test frame with original pose $(\mathbf{R}, \mathbf{T})$ and ground
truth $I$, and for an evaluation elevation $e$ (and, optionally, one of the
four cardinal azimuths, left, right, up, or down), RTVC proceeds as
follows.
\begin{enumerate}
\item \textbf{Render.} The sideview pose $(\mathbf{R}_{a,e},
\mathbf{T}_{a,e})$ is synthesized as in Section~\ref{sec:method:gcsr}, and
the trained model is rendered at this pose to obtain a color image
$\tilde{I}$ and depth map $\tilde{D}$.
\item \textbf{Warp back.} $(\tilde{I}, \tilde{D})$ is forward-warped from
the sideview pose to the original pose using a depth-based, $z$-buffered
point splat: each valid sideview pixel is unprojected to a 3D point using
$\tilde{D}$ and the known sideview camera intrinsics and pose, reprojected
into the original camera, and rasterized with nearest-wins $z$-buffering to
resolve duplicate or occluded projections. This produces a warped image
$\hat{W}$ and a validity mask $M$ indicating which pixels of the original
frame the sideview render's geometry was able to re-explain.
\item \textbf{Score.} $\hat{W}$ is compared against the ground truth $I$
using masked PSNR and masked SSIM, restricted to pixels valid under both
$M$ and the surgical-tool mask, and, for SSIM, to windows whose valid
coverage exceeds a minimum threshold.
\end{enumerate}

\paragraph{Ground-truth-free property.}
No photograph at the sideview pose is required at any point in this
procedure: the only ground truth used is $I$, which was already captured
at the original pose. The resulting comparison is nonetheless a genuine
photometric evaluation, because any error in the reconstructed geometry at
the sideview pose, such as incorrect depth or missing or hallucinated
structure, displaces content relative to the true pixels of $I$ once
warped back, degrading RTVC's PSNR and SSIM even when the sideview render
$\tilde{I}$ appears plausible in isolation and cannot itself be validated
against a captured photograph.

\paragraph{Statistical evaluation.}
For a given elevation $e$, the up to four cardinal-azimuth measurements are
averaged per frame prior to any statistical comparison, as they share the
same underlying reconstruction at that time step and do not constitute
independent samples; treating them as independent would pseudo-replicate
the data and overstate statistical significance. This yields a single
paired sample per frame per model, which is compared between two models
using a paired $t$-test across frames. The special case $e = 0^\circ$,
corresponding to no sideview offset and no warp, reduces to the
conventional test-view PSNR and SSIM, and is reported alongside the RTVC
elevations as a reference point.
